% \subsection{Reinforcement Learning for Dexterous Grasping}
% Reinforcement learning (RL) has emerged as a powerful paradigm for acquiring
% complex manipulation skills in high-degree-of-freedom (DoF) systems such as
% dexterous robotic hands. Recent research has relied heavily on model-free RL
% algorithms, including DDPG, SAC, and PPO, which have demonstrated strong
% performance in learning a variety of skills across a wide range of robotic
% setups.

% GraspXL reformulates the RL problem as an objective-conditioned
% motion generation task. The policy is trained to simultaneously
% produce stable grasps while satisfying user-specified motion
% objectives, such as graspable regions defined via signed
% distance fields (SDFs), approach directions, wrist rotations,
% and target hand poses. This formulation enables multi-objective
% policy learning within a unified reinforcement learning
% framework.

% DextrAH-RGB\cite{singh2024dextrah} adopts a geometric fabrics-based reinforcement
% learning framework, in which the policy operates within a
% hand-designed fabric-induced action manifold. This structured
% representation effectively reduces the action-space dimension
% and significantly simplifies reward design during training.

\textbf{Skill Distillation.} Skill distillation transfers behavior from a high-capacity teacher to a compact student and has been widely adopted
in both computer vision and robotics. In robotics, the teacher can leverage privileged state and simulator access during training,
while the student is distilled to operate from on-board observations (e.g., RGB). DexTrAH-RGB~\cite{singh2024dextrah} 
distills a privileged RL teacher into an RGB-based policy via DAgger-based imitation learning, while 
UniGraspTransformer~\cite{wang2025unigrasptransformer} distills offline trajectories from object-specific teachers 
into a unified Transformer policy.

\textbf{Adaptive Intervention in Interactive Imitation Learning.} Many imitation learning methods use fixed offline datasets for supervision~\cite{zhang2025contactdexnet,zhang2025omnidexvlg}. In sequential decision-making, however, distribution shift induces covariate shift and compounding errors. Interactive imitation learning~\cite{welte2025interactive} mitigates these issues by collecting online supervision along trajectories visited by the learner.
DAgger~\cite{ross2011reduction} addresses covariate shift by iteratively relabeling learner-visited states with an expert, while SafeDAgger~\cite{zhang2016query} adds teacher intervention in high-risk states to reduce unsafe exploration.

Later methods introduce adaptive intervention mechanisms based on uncertainty-driven rules~\cite{menda2019ensembledagger}, human gating~\cite{kelly2019hgdagger}, or learned intervention criteria~\cite{cai2025robot}. However, they have primarily been evaluated in relatively simple control or autonomous-driving settings. Fully autonomous expert relabeling during online distillation remains rare in robotic manipulation,
especially for dexterous, contact-rich tasks.

\textbf{Vision-Language-Guided Robot Learning.}
Visual and semantic information is widely recognized as a strong 
cue for embodied agents to interpret task specifications and reason 
about unstructured environments~\cite{firoozi2025foundation,han2026multimodal}.
One line of work uses VLMs to generate demonstrations or supervisory signals
for offline policy construction, such as
Manipulate-Anything~\cite{duan2025manipulate}.
Complementary approaches, including RACER~\cite{dai2025racer} and
FOREWARN~\cite{wu2025foresight}, employ VLMs to guide risk prediction,
provide recovery guidance for unsafe situations, or govern action selection
during online policy learning or real-time execution.
